\section{Introduction}
\label{sec:intro}

Quantum and quantum-inspired annealers, such as D-Wave's quantum processors, Fujitsu's Digital
Annealer~\citep{aramon2019digital} and GPU-based annealers, all accept the same input: a quadratic unconstrained
binary optimization (QUBO) problem. Many NP-hard problems have known QUBO formulations~\citep{lucas2014ising,
glover2019tutorial}, which makes QUBO an attractive common interface between real decision problems and new
hardware. In practice, however, a usable QUBO solution of a real problem requires a chain of expert work:
\begin{enumerate}
  \item pinning down what is actually being decided, under which rules (often laws or standards);
  \item finding real data, cleaning it and bringing it to a size the hardware can handle;
  \item writing a formulation that captures the confirmed objective and constraints;
  \item choosing penalty weights, encodings of inequality constraints and solver parameters, which strongly
        affect solution quality on hardware~\citep{verma2020penalty,ayodele2022penalty,pelofske2025comparing};
  \item explaining the result to the person who asked.
\end{enumerate}

Large language models (LLMs) have recently been used to automate step~3. QuantumQUBO
Agent~\citep{mondal2026quantumqubo} turns a natural-language description into a QUBO with a multi-agent pipeline
and verifies it by brute force on small test cases; related work does the same for mixed-integer
programming~\citep{ramamonjison2023nl4opt,ahmaditeshnizi2024optimus,xiao2024chainofexperts}. These systems stop at
the formulation, evaluate on toy-sized instances (in the QuantumQUBO benchmark, the median test case has five
binary variables), use large proprietary models, and list execution on hardware as future work.

Asking a frontier chat model the same question shows a different gap. It will write an objective and
constraints from one sentence and point to public datasets, but the rest of the work stays with the user: finding
and downloading the data, cleaning and joining it, running the solver, plotting and tuning. Our goal is to
automate exactly that remaining work, for real problems with real data, using a model that runs on a local GPU.

\paragraph{This paper.}
We build and study a human-in-the-loop system that takes a one-sentence problem (``place 20 EV chargers in
Taipei'', ``schedule 12 nurses for four weeks under Taiwan's Labor Standards Act'', ``plan a day of sightseeing in
Paris'') through data preparation, modelling, annealing and delivery. All agents run on one locally served open
model. People confirm the specification, approve data assumptions and review the result. The operational goal is
to complete this workflow or identify a concrete unresolved gap; it does not require a certified optimum for
every application. Poor results may trigger further solver search or a return to data and model repair,
depending on the evidence.

\paragraph{Research question.}
How can a locally served LLM, supported by specialized agent roles and executable feedback, turn a confirmed
request into a correctly prepared instance, a requirement-consistent model and a usable solution? Our design
separates these three obligations: feasibility of the generated model alone cannot establish the first two.
Multi-agent coordination and retrieval are mechanisms to evaluate, rather than benefits assumed from the
number of roles or the presence of a knowledge library.

\paragraph{Contributions.}
\begin{enumerate}
  \item \textbf{An end-to-end, local-LLM system from a sentence to an annealer (\cref{sec:system}).}
        Clarification with sourced web lookup; autonomous formulation under problem-independent structural
        checks; real-data acquisition and preprocessing with a tested toolkit and a provenance record;
        an encoding library with programmatic guardrails; one interface to four backend families (simulated
        annealing, a GPU annealer, the Fujitsu Digital Annealer and D-Wave QPU or Leap Hybrid), of which
        simulated annealing, the GPU annealer and Leap Hybrid were used in our experiments; and generated figures and
        explanations. Bounded recovery routes failures to data preparation, modelling or solving, retains
        versioned artifacts, and reports delivery status against a feasible conventional baseline.
  \item \textbf{A data model for ``real problems'' (\cref{sec:setting}).} Each input of a model is either
        public data (found, never invented), a user setting (asked, with defaults) or private data (uploaded,
        or explicitly marked mock data). This makes ``uses real data'' a checkable property of a run instead of
        a claim, and it decides what the system automates and what it asks. The current implementation declares
        entity sets, indexed parameters, units and derivations explicitly; a separate Data Checker can compare
        counts and totals with computations on the raw tables.
  \item \textbf{Problem-level verification of LLM-written QUBOs (\cref{sec:verification}).} We show that
        bitstring-based verification with LLM-derived reference solutions misjudges models in both directions,
        and that the reference answers of a public benchmark can themselves be inconsistent. We replace it with
        a check that decodes every ground state and evaluates it with an independent checker anchored on the
        human-stated optimum, and we quantify run-to-run instability that single-run evaluations hide. For
        larger real-data instances, a separate Requirement Checker samples decisions against the confirmed
        request and checks delivered decisions; this sampled check does not certify all model states.
  \item \textbf{Failure modes of a small local model as an autonomous modeller, and mechanisms that address
        them (\cref{sec:system,sec:experiments}).} We document recurring failures (empty objectives, constraint
        penalties folded into the objective, silent over-simplification into a sortable problem, dropped hard
        constraints, exhausted reasoning budgets, prompt overflow on large models) and the generic mechanisms
        that address them: task splitting, repair instead of rewrite, constraint families and prompt budgets.
        The current system also implements bounded data-to-model repair and optional retrieval of general
        modelling patterns; their end-to-end effects require the ablations in \cref{sec:autonomousprotocol}.
  \item \textbf{An empirical characterization (\cref{sec:experiments}).} A 12-category end-to-end suite, the
        scaling of a GPU annealer on routing up to 3{,}600 variables, and a budget-matched comparison of
        LLM-driven tuning with rule-based, random and Bayesian-optimization baselines, which shows that penalty
        tuning benefits depend on the instance and backend (\cref{sec:penaltysummary}; details in \cref{app:details}).
\end{enumerate}

The current autonomous acceptance suite targets Taipei EV-charger coverage, a YouBike inspection tour and AED
inspection grouping. Its evaluation judge is separate from the runtime agents. Development runs have exposed
both successful deliveries and upstream failures; they are not pooled into a success rate because code and
prompts changed between runs. The earlier synthetic suite and specified-source pilots remain separate evidence
(\cref{sec:autonomousprotocol,sec:limitations}).
